% FREUID Challenge 2026 — Technical Report, team junesdata
% Compile: pdflatex freuid_technical_report.tex && bibtex freuid_technical_report && pdflatex x2

\documentclass[11pt,a4paper]{article}

\usepackage[T1]{fontenc}
\usepackage[utf8]{inputenc}
\usepackage{lmodern}
\usepackage{microtype}
\usepackage{geometry}
\usepackage{graphicx}
\usepackage{booktabs}
\usepackage{amsmath}
\usepackage{hyperref}
\usepackage{xcolor}
\usepackage{enumitem}

\geometry{margin=2.5cm}
\hypersetup{
  colorlinks=true,
  linkcolor=blue!60!black,
  citecolor=blue!60!black,
  urlcolor=blue!60!black,
}

\newcommand{\teamname}{junesdata}
\newcommand{\kaggleteam}{junesdata}
\newcommand{\contactemail}{tinkode6@gmail.com}
\newcommand{\reportdate}{July 15, 2026}

\title{%
  \textbf{The FREUID Challenge 2026}\\[0.4em]
  \large Technical Report --- \teamname
}
\author{%
  Necula George-Andrei\\
  \small Independent\\
  \texttt{\contactemail}
}
\date{Report date: \reportdate}

\begin{document}
\maketitle

\begin{abstract}
We address binary fraud detection on identity-document images. Our key observation is
that fraud signal in FREUID is dominantly \emph{local}: manipulations live at the level
of pixels and small regions, while document-level semantics occur in both classes.
Consequently we avoid full-image downscaling entirely and train binary patch
classifiers (EfficientNetV2-S and EfficientNetV2-M, ImageNet-21k pre-trained) on random
$512\times512$ crops taken at native resolution, with a recapture-simulation
augmentation (resampling, moir\'e-like patterns, sensor noise, JPEG recompression)
targeting the capture-pipeline shift emphasized by the test set. At inference each
image is covered by an overlapping grid of $512\times512$ crops and scored by the mean
of its top-2 crop logits; the final score is the mean of per-model ranks, a scale-free
combination well matched to the ranking-based FREUID Score (lower is better). Training
uses only the official FREUID training set plus iterative pseudo-labeling on the
official public test images. Full-image baselines scored $0.20$--$0.25$ on the public
leaderboard; the patch pipeline with pseudo-labeling and the two-model rank ensemble
reached $\mathbf{0.00226}$. Code (MIT), frozen weights, and a no-network Docker
container reproducing the selected submissions are public.
\end{abstract}

\noindent\textbf{Competition:} The FREUID Challenge 2026 (IJCAI-ECAI 2026)\\
\textbf{Kaggle team:} \kaggleteam\\
\textbf{Code:} \url{https://github.com/4kaws/freuid-challenge-2026} (commit \texttt{[FINAL-SHA]})\\

\section{Introduction}
\label{sec:intro}
The FREUID Challenge~\cite{freuid2026} asks for a calibrated fraud score for
identity-document images under a realistic threat model: physical manipulations,
GenAI-driven edits, and print-and-capture forgeries that suppress fragile digital
artifacts. The evaluation emphasizes generalization: the private test set contains
document types absent from both training data and public test data, and places
stronger weight on recaptured (non-synthetic) examples.

Early experiments showed that models trained on downscaled full images learn
shortcuts that do not survive the distribution shift between the training generation
batch and the test batch (public-leaderboard FREUID $\approx 0.20$--$0.25$ despite
saturated in-distribution validation). Manual inspection indicated that semantic
anomalies (implausible fields, mismatched attributes) occur in \emph{both} classes,
i.e.\ they are generator noise rather than fraud signal, while true fraud is local
pixel-level manipulation. This motivated a native-resolution patch approach.

\section{Method}
\label{sec:method}

\subsection{Overview}
The pipeline is: native-resolution $512\times512$ patch classifier(s)
$\rightarrow$ overlapping grid-crop inference with top-$k$ logit pooling per image
$\rightarrow$ cross-model rank averaging $\rightarrow$ submission scores in $[0,1]$.

\subsection{Model architecture}
Two timm~\cite{timm} backbones fine-tuned as binary classifiers
(\texttt{num\_classes=1}):
\begin{itemize}[noitemsep]
  \item \texttt{tf\_efficientnetv2\_s.in21k\_ft\_in1k} ($\approx$21M parameters);
  \item \texttt{tf\_efficientnetv2\_m.in21k\_ft\_in1k} ($\approx$54M parameters).
\end{itemize}
Inputs are $512\times512$ RGB crops normalized with ImageNet statistics. No resizing
is applied anywhere: crops are taken from images at their original resolution
(images smaller than 512 px on a side are zero-padded).

\subsection{Training objective and procedure}
Binary cross-entropy with logits; AdamW (lr $3\times10^{-4}$ for V2-S with batch 48,
$2\times10^{-4}$ for V2-M with batch 24; weight decay $10^{-2}$); OneCycle schedule;
4 epochs; mixed precision. One random $512\times512$ crop per image per epoch with
horizontal/vertical flips. With probability $0.3$ a \emph{recapture simulation} is
applied to the crop: random down/up-resampling ($0.45$--$0.9\times$, mixed
interpolation), moir\'e-like sinusoidal luminance patterns, Gaussian sensor noise,
gamma/color drift, and JPEG recompression (quality 45--92). This augmentation targets
the print-and-capture emphasis of the test distribution.

\textbf{Iterative pseudo-labeling.} The public test images (unlabeled competition
data) are incorporated by self-training: after each round, images with fraud
probability $>0.98$ are pseudo-labeled fraud, and the 1{,}150 lowest-ranked images are
pseudo-labeled genuine (rank-based rather than threshold-based, because confident
genuine calibration collapses after the first round). Five rounds were run; each round
retrains from ImageNet initialization on train + current pseudo-labels. Public
leaderboard progression across rounds: $0.0587 \rightarrow 0.0350 \rightarrow 0.0224
\rightarrow 0.0038 \rightarrow 0.0035$ (with inference improvements applied along the
way, see Section~\ref{sec:inference}).

\subsection{Score calibration}
The FREUID Score is computed from the DET curve and an operating point
(APCER@1\%BPCER), both invariant to strictly monotone transforms of the scores.
We therefore use rank-normalized scores: for each model, per-image aggregated logits
are converted to ranks in $[0,1]$ over the full test set; the submission score is the
(weighted) mean of per-model ranks, clipped to $[10^{-6}, 1-10^{-6}]$.

\section{Data}
\label{sec:data}

\subsection{FREUID training set}
All 69{,}352 training images across the five document types were used. For model
selection during development we used leave-one-type-out folds (train on four types,
validate on the held-out type), which correlates with the public leaderboard far
better than random splits (full-image models: held-out-type FREUID $\approx 0.07$--$0.22$
vs.\ in-distribution $\approx 0$).

\subsection{External data and pre-trained models}
\emph{No external data beyond FREUID and publicly licensed pre-trained weights.}
Pre-trained backbones are the public timm ImageNet-21k/1k checkpoints
(Apache-2.0)~\cite{timm}.

\subsection{Validation strategy}
Development signals: (i) leave-one-type-out validation on the training types;
(ii) the public leaderboard, used as the primary signal for the pseudo-labeling
rounds (the public test set is explicitly designated as a development signal by the
organizers). Private test labels were never available; private test images were not
used for any training or tuning --- all model weights were finalized and archived
(timestamped Kaggle dataset \texttt{junesdata/freuid-ckpts}) \emph{before} the private
test release on July 13, 2026, 07:02 UTC.

\section{Inference}
\label{sec:inference}
Each test image is covered by an overlapping $R\times C$ grid of $512\times512$ crops
at native resolution (grid offsets evenly spaced; images smaller than the crop are
zero-padded). Per model, the image score is the mean of the top-2 crop logits ---
max-style pooling is essential because fraud is local (mean pooling degrades the
public score by $\sim3\times$). Grid density improved the public score up to
$7{\times}9=63$ crops, with saturation beyond. The submitted configurations balance
accuracy against the 6-hour A100 verification budget.

Two final picks are produced from one Docker image and one frozen weight set,
differing only in a documented environment flag (inference orchestration only):
\begin{itemize}[noitemsep]
  \item \textbf{Pick 1} \texttt{VARIANT=ens2}: rank-mean of EfficientNetV2-S (round 5)
        and EfficientNetV2-M (round 5), grid \texttt{[FINAL-GRID-1]}.
  \item \textbf{Pick 2} \texttt{VARIANT=hedge}: adds the round-1 pseudo-label model and
        the no-pseudo-label base model at half weight, grid \texttt{[FINAL-GRID-2]} ---
        an out-of-distribution hedge that reduces specialization to public-test fraud
        patterns.
\end{itemize}
Commands and output checksums: see Table~\ref{tab:picks} and the repository README.

\begin{table}[h]
  \centering
  \caption{Final picks: submission $\leftrightarrow$ command $\leftrightarrow$ checksum.}
  \label{tab:picks}
  \small
  \begin{tabular}{lll}
    \toprule
    Pick & Command & \texttt{sha256(submission.csv)} \\
    \midrule
    1 (\texttt{ens2})  & \texttt{docker run --network none ... -e VARIANT=ens2}  & \texttt{[CHECKSUM-1]} \\
    2 (\texttt{hedge}) & \texttt{docker run --network none ... -e VARIANT=hedge} & \texttt{[CHECKSUM-2]} \\
    \bottomrule
  \end{tabular}
\end{table}

\section{Results}
\label{sec:results}

\begin{table}[h]
  \centering
  \caption{Public-leaderboard progression (FREUID Score; lower is better).}
  \label{tab:results}
  \begin{tabular}{lcc}
    \toprule
    Configuration & Public FREUID $\downarrow$ & Notes \\
    \midrule
    Full-image EfficientNetV2-S @384/@768        & 0.250 / 0.202 & shortcut learning \\
    Patch 512, $2{\times}3$ grid, top-2          & 0.0737 & \\
    Patch 512, $3{\times}4$ grid, top-2          & 0.0587 & \\
    + pseudo-labels (round 1)                    & 0.0350 & \\
    + denser grid $5{\times}6$                   & 0.0276 & \\
    + pseudo-labels (round 3)                    & 0.0224 & \\
    + round 4, $7{\times}9$ grid                 & 0.0038 & \\
    + round 5, $7{\times}9$ grid                 & 0.0035 & \\
    + two-model rank ensemble (V2-S + V2-M)      & \textbf{0.00226} & \\
    Final Pick 1 (\texttt{ens2})                 & [TBD]  & Selected final submission \\
    Final Pick 2 (\texttt{hedge})                & [TBD]  & OOD-robust variant \\
    \bottomrule
  \end{tabular}
\end{table}

\section{Reproducibility}
\label{sec:repro}
\begin{itemize}[noitemsep]
  \item \textbf{Repository:} \url{https://github.com/4kaws/freuid-challenge-2026},
        commit \texttt{[FINAL-SHA]} (MIT license).
  \item \textbf{Docker:} \texttt{docker build -f docker/Dockerfile -t freuid-solution .};
        weights are inside the repository (Git LFS) and are copied into the image ---
        no runtime downloads.
  \item \textbf{Dependencies:} PyTorch 2.3.1 + CUDA 12.1 base image; timm 1.0.11;
        OpenCV (headless); pinned in \texttt{docker/requirements.txt}.
  \item \textbf{Runtime:} single A100 40GB; measured \texttt{[RUNTIME]} for the full
        142{,}818-image test set at the submitted configuration (within the 6-hour
        verification budget). CPU fallback exists but is not practical at scale.
  \item \textbf{Randomness:} inference is deterministic up to GPU non-associativity;
        no test-time augmentation with random components is used.
  \item \textbf{Network:} the container runs with \texttt{--network none}; verified
        locally before submission.
  \item \textbf{Code freeze:} all model weights used in the final picks were archived
        in the version history of the Kaggle dataset \texttt{junesdata/freuid-ckpts}
        before the private test release (last weight version: July 13, 2026, 04:56 UTC).
\end{itemize}

\section{Team and contributions}
Necula George-Andrei --- modeling, data analysis, engineering, and writing.

\section*{Acknowledgments}
We thank the organizers and the Microblink Fraud Lab for a well-designed competition.

\bibliographystyle{plain}
\bibliography{references}

\appendix
\section{Hyperparameters}
\label{app:hparams}
\begin{table}[h]
  \centering
  \small
  \begin{tabular}{ll}
    \toprule
    Hyperparameter & Value \\
    \midrule
    Crop size & $512\times512$ (native resolution, no resizing) \\
    Batch size & 48 (V2-S) / 24 (V2-M) \\
    Learning rate & $3\times10^{-4}$ (V2-S) / $2\times10^{-4}$ (V2-M), OneCycle \\
    Weight decay & $10^{-2}$ \\
    Epochs & 4 per round \\
    Recapture augmentation prob. & 0.3 \\
    Pseudo-label thresholds & fraud $p>0.98$; genuine = 1150 lowest-ranked \\
    Inference pooling & mean of top-2 crop logits \\
    Ensemble & rank mean over models \\
    Seed & 42 (training); inference deterministic \\
    \bottomrule
  \end{tabular}
\end{table}

\section{Negative results}
For completeness: per-type rank normalization of a full-image model (0.83), mean
pooling over crops (0.21), ensembling leave-one-type-out fold models into the test
ensemble (dilutes: 0.075 vs.\ 0.035 baseline), agreement-gated pseudo-label round 2
(0.043), a 6th blind pseudo-label round without leaderboard feedback (0.0055 vs.\
0.0035), and denser grids beyond $7{\times}9$ (flat: 0.0036 at $9{\times}12$).

\end{document}
